\RequirePackage{iftex}
\ifPDFTeX
\documentclass[a4paper,12pt]{article}
\usepackage[margin=22mm]{geometry}
\usepackage[T1]{fontenc}
\usepackage[utf8]{inputenc}
\usepackage{graphicx}
\usepackage{CJKutf8}
\else
\documentclass[a4paper,12pt]{jsarticle}
\usepackage[dvipdfmx,margin=22mm]{geometry}
\usepackage[dvipdfmx]{graphicx}
\fi
\usepackage{amsmath,amssymb,amsthm}
\usepackage{enumitem}
\usepackage{tikz}
\usetikzlibrary{arrows.meta,calc}

\setlength{\parindent}{0pt}
\setlength{\parskip}{0.35em}
\setlength{\fboxsep}{8pt}
\setlist[itemize]{leftmargin=2em,itemsep=0.25em,topsep=0.35em}
\setlist[enumerate]{leftmargin=2.2em,itemsep=0.45em,topsep=0.35em}

\newcommand{\feedbackqa}[2]{%
\item[]
\begin{tabular}[t]{@{}p{1.8em}@{\hspace{0.4em}}p{\dimexpr\linewidth-2.2em\relax}@{}}
Q. & #1 \\
A. & #2
\end{tabular}
\par\smallskip\hrulefill
}

\newcommand{\feedbacknote}[1]{%
\item[] #1\par\smallskip\hrulefill
}

\newcommand{\pointbox}[1]{%
\begin{center}
\fbox{%
\begin{minipage}{0.91\linewidth}
#1
\end{minipage}}
\end{center}
}

\begin{document}
\ifPDFTeX
\begin{CJK}{UTF8}{min}
\fi

\theoremstyle{definition}
\newtheorem*{definition}{定義}
\newtheorem*{theorem}{定理}
\newtheorem*{example}{例}
\newtheorem*{remark}{注意}
\renewcommand{\proofname}{\normalfont 証明}

\begin{center}
{\LARGE 数学1A 第14回レジュメ}\\[0.4em]
\end{center}

\section*{フィードバック}

\begin{itemize}[leftmargin=1.5em,itemsep=0.7em,topsep=0.3em]
\feedbackqa{$G(x,y)=x^3-y^2=0$ のもとで、$F(x,y)=x$ の最小値は？}{非常に良い問題ですね。ラグランジュ関数を\[L(x,y,\lambda)=x-\lambda(x^3-y^2)\]とおきます。通常のラグランジュの未定乗数法では、\[L_x=1-3\lambda x^2=0,\qquad L_y=2\lambda y=0,\qquad x^3-y^2=0\]を連立して解きます。第1式から $\lambda\ne0$ なので、第2式より $y=0$ です。すると制約条件から $x=0$ となりますが、第1式は $1=0$ となり矛盾します。つまり、$\nabla G\ne0$ である通常の候補点はありません。ここで $(0,0)$ では $\nabla G(0,0)=(0,0)$ となり、ラグランジュの未定乗数法の正則性の仮定が成り立たないため、この点を別に調べる必要があります。制約 $x^3-y^2=0$ は $y^2=x^3$ と書けるので、制約を満たす点では $x\geq0$ です。また $(0,0)$ は制約を満たし、$F(0,0)=0$ です。したがって、\[{\text{最小値は }0\text{、それを与える点は }(0,0)}\]となります。}
\feedbackqa{夏休みに先取りをしたほうがよいか？}{時間をかけて、先取りすべきかはその人次第ですが、少しの時間（10分くらい）でも何をやるのかを見ておくだけでも、理解の進度がだいぶ変わると思います。}
\feedbackqa{1限の授業だったが、毎回出席できた。}{お疲れさまでした。1限は大変ですが、継続した経験はこれからにつながると思います。}
\feedbackqa{ラグランジュの $\lambda$ は常に $\lambda\ne0$ か？}{常に $0$ でないとは限りません。ラグランジュの未定乗数法では、候補点で\[\nabla f=\lambda\nabla g\]が成り立ちます。$\lambda=0$ のときは $\nabla f=0$ となるだけで、矛盾ではありません。例えば、制約 $x^2+y^2=1$ のもとで $f(x,y)=x^2$ を最小にすることを考えます。$(0,1)$ と $(0,-1)$ では $\nabla f=(2x,0)=(0,0)$ なので、上の式は $\lambda=0$ で成り立ちます。実際、この2点で最小値 $0$ をとります。したがって、途中で $\lambda$ で割るときは、先に $\lambda=0$ の場合を除外できるか確認する必要があります。今回の解答のように $3-2\lambda x=0$ という式がある場合は、$\lambda=0$ なら $3=0$ となってしまうため、$\lambda\ne0$ と分かります。}
\feedbackqa{教科書の問題が初見で解けない。普段どのように教科書を読んでいるか？}{解けずに悩むこと自体が大切です。今は解けなくても、少し時間を置くと解けるようになることはよくあります。教科書を読むときは、「なぜそうなるのか」「これは本当に正しいのか」と批判的に読むのがよいといわれています。}
\feedbackqa{教科書の問題はどの程度解くべきか？}{できればすべて解くのがよいと思います。時間がなければ、基礎問題を自分で選んで解きましょう。チュートリアルアワーで相談してもよいでしょう。}
\feedbackqa{黒板の端に書かれると見えない。}{ご指摘ありがとうございます。気を付けます。}
\feedbackqa{やり方が分かってもゴールが見えず、何をしているのか分からない。}{目標を先に理解することは大切です。あらかじめ誰かに聞いたり、自分で調べたりすることを勧めます。}
\feedbackqa{小テストの問題を LMS にアップしてほしい。}{小テストの問題は毎回レジュメに載せているので、そちらを確認してください。}
\feedbackqa{ラグランジュはなぜ変数を増やそうと思ったのか？}{物理（特に熱力学）では、補助変数という新しい変数を導入して式を書き換えることがよくあったので、そこから来ているような気がします。ただ、天才が考えることはよくわかりません。}
\end{itemize}

\section*{小テスト}

\begin{enumerate}
\item 曲線
\[
x^2+xy+y^2=7
\]
について、次の問いに答えよ。
\begin{enumerate}[label=(\alph*)]
\item $\dfrac{dy}{dx}$ を $x,y$ を用いて表せ。
\item 点 $(1,2)$ における接線の方程式を求めよ。
\end{enumerate}

\item 制約条件 $x^2+y^2=1$ のもとで、関数
\[
f(x,y)=3x+4y
\]
の最大値と、それを与える点をラグランジュの未定乗数法を用いて求めよ。
\end{enumerate}

\section*{解答}

\subsection*{1. 陰関数の微分}

与えられた式
\[
x^2+xy+y^2=7
\]
を $x$ で微分すると、
\[
2x+\left(x\frac{dy}{dx}+y\right)+2y\frac{dy}{dx}=0
\]
となる。したがって、
\[
\left(x+2y\right)\frac{dy}{dx}=-(2x+y),
\qquad
\frac{dy}{dx}=-\frac{2x+y}{x+2y}
\]
を得る。

点 $(1,2)$ では
\[
\frac{dy}{dx}=-\frac{2\cdot1+2}{1+2\cdot2}=-\frac45
\]
である。よって、この点における接線は
\[
y-2=-\frac45(x-1)
\]
である。

\begin{center}
\begin{tikzpicture}[x=0.9cm,y=0.9cm]
\draw[->] (-4.2,0) -- (4.2,0) node[right] {$x$};
\draw[->] (0,-3.8) -- (0,3.8) node[above] {$y$};
\draw[blue!65!black,thick,samples=120,domain=0:360,smooth,variable=\t]
  plot ({(sqrt(14/3)*cos(\t)+sqrt(14)*sin(\t))/sqrt(2)},
        {(sqrt(14/3)*cos(\t)-sqrt(14)*sin(\t))/sqrt(2)});
\draw[red!70!black,thick,domain=-1.2:4.2,samples=2]
  plot (\x,{-0.8*\x+2.8});
\fill[red!70!black] (1,2) circle (2pt) node[above right] {$(1,2)$};
\node[blue!65!black] at (-2.3,-3.1) {$x^2+xy+y^2=7$};
\node[red!70!black] at (3.0,0.1) {接線};
\end{tikzpicture}
\end{center}

青が曲線 $x^2+xy+y^2=7$、赤が点 $(1,2)$ における接線である。

\subsection*{2. ラグランジュ未定乗数法（最大値）}

制約条件を $x^2+y^2=1$ とする。ラグランジュ関数を
\[
\Phi(x,y,\lambda):=3x+4y-\lambda(x^2+y^2-1)
\]
と定義する。極値候補では各変数に関する偏微分がすべて $0$ となるので、
\[
\Phi_x=3-2\lambda x=0,\qquad
\Phi_y=4-2\lambda y=0,\qquad
\Phi_\lambda=-(x^2+y^2-1)=0
\]
を連立して解けばよい。最初の二式から
\[
x=\frac{3}{2\lambda},\qquad y=\frac{2}{\lambda}
\]
であり、これを制約条件に代入すると
\[
\lambda=\pm\frac52
\]
となる。したがって、極値候補は
\[
(x,y)=\left(\frac35,\frac45\right),\qquad
\left(-\frac35,-\frac45\right)
\]
である。それぞれでの関数値は $5,-5$ なので、最大値を与える点と最大値は
\[
{(x,y)=\left(\frac35,\frac45\right)},
\qquad
{\text{最大値は }5}
\]
である。

\ifPDFTeX
\end{CJK}
\fi
\end{document}
